\section{Short curves and small eigenvalues}\label{sec:geometry}

The analytic estimates enter geometry through ratios of volumes of cut
surfaces, by Mirzakhani's integration formula~\cite{MirInv}. We first apply the
adjacent-volume limit to Shen--Wu's vanishing-gap argument~\cite{SW}
(Subsection~\ref{sec:spectral}). We then determine the short-curve process
and the systole above the scale $\sqrt g$
(Subsection~\ref{sec:pants-spectrum}) and deduce many small eigenvalues
(Subsection~\ref{sec:pants-eigenvalues}). The last two subsections state
the limit of the tight length spectrum at every cusp density and of the
full primitive spectrum at the critical scale $n\asymp\sqrt g$; their
proofs are given in Appendices~\ref{app:tight} and~\ref{app:critical}. The
critical result uses volume bounds of Mirzakhani--Zograf and Shen--Wu
and only the zero-density case of the volume-ratio limits, not the
uniform expansion.

\subsection{The first positive eigenvalue}\label{sec:spectral}

The simplest spectral consequence uses only the discrete cusp ratio
$D_{g,n-1}$. A short curve enclosing two cusps bounds a pair of pants
of fixed area $2\pi$, joined to the rest of the surface through a short
boundary. This gives a small Cheeger ratio and hence a small Rayleigh
quotient. Shen--Wu's second-moment argument makes such a curve probable.
Shen and Wu asked for the adjacent-volume limit needed to extend
this argument beyond $n=o(g)$~\cite[Question~3]{SW}. Theorem~\ref{thm:volume}
provides that limit along every sequence with $g\to\infty$.

\begin{proof}[Proof of Corollary~\ref{cor:spectral}]
Put $n=n(g)$ and
\[
 L_g=\left(\frac{\sqrt g}{n}\right)^{1/2}.
\]
Let $N_g(X)$ count the simple closed geodesics of length at most $L_g$ whose complement has components of topological types $(0,3)$ and $(g,n-1)$. Proposition~35 of~\cite{SW}, in the notation of its arXiv version~4, states under the sole assumption $n/\sqrt g\to\infty$ that
\[
 \liminf_{g\to\infty}\mathbb P^{\mathrm{WP}}_{g,n}\{N_g\ge1\}
 \ge \liminf_{g\to\infty}
 \frac{V_{g,n-1}^{2}}{V_{g,n}V_{g,n-2}}.
\]
The ratio on the right is $D_{g,n-1}$, which tends to one by Theorem~\ref{thm:volume}. Thus $\mathbb P^{\mathrm{WP}}_{g,n}\{N_g\ge1\}\to1$. The same conclusion also follows within this paper: $N_g$ is the count $N_{L_g}$ of Subsection~\ref{sec:pants-spectrum}, with $q_g$ as in \eqref{eq:PS1}, and since $L_g\to0$ and $q_gL_g^2\asymp n\sqrt g/(g+n)\to\infty$, Proposition~\ref{prop:pants-concentration} gives $N_g/(q_gL_g^2)\to1$ in probability. Its proof is the same second-moment argument, with $D_{g,n-1}\to1$ as the volume input.

Write $\mathfrak h(X)$ for the Cheeger constant. On the event $N_g\ge1$, a counted geodesic bounds a hyperbolic pair of pants of area $2\pi$; the other component has area $2\pi(2g+n-3)\ge2\pi$ for all sufficiently large $g$. This separating curve is an admissible Cheeger cut, so
\[
 \mathfrak h(X)\le \frac{L_g}{2\pi}.
\]
Let $\mathcal R(X)$ be the infimum of the Rayleigh quotient over nonzero $H^1(X)$ functions orthogonal to constants. Buser's inequality~\cite[Theorem~7.1]{Buser}, in the form recorded in~\cite[Theorem~11]{SW} for complete noncompact hyperbolic surfaces, gives $\mathcal R(X)\le c\mathfrak h(X)$ with an absolute constant $c$. Since $L_g\to0$, this upper bound is eventually smaller than both $\varepsilon$ and $1/4$. The essential spectrum starts at $1/4$; the min--max principle, as recalled in~\cite[Theorem~7]{SW}, therefore yields a positive discrete eigenvalue $\lambda_1(X)=\mathcal R(X)<\varepsilon$ on $N_g\ge1$, for all sufficiently large $g$. This proves the asserted probability limit.
\end{proof}

The use of $D_{g,n-1}$ above is the only change to the final step of Shen--Wu's proof: under their condition $n=o(g)$ this limit follows from \cite[Lemma~5.1]{MZ}, as they note. The short-curve estimate, the Cheeger cut, and the passage to a discrete eigenvalue are their geometric and spectral inputs.

\subsection{Short-curve processes and the systole}
\label{sec:pants-spectrum}

We now determine the limiting process of short curves and the systole.
The moment estimates proved here will also be used in
Subsection~\ref{sec:pants-eigenvalues} to count many disjoint short pairs of pants.
Throughout this and the next subsection,
$g\to\infty$, $n/\sqrt g\to\infty$, and $X$ has law
$\mathbb P^{\mathrm{WP}}_{g,n}$. Let $\mathcal P_2(X)$ be the set of
unoriented simple closed geodesics bounding a pair of pants with two
cusps. Define the deterministic intensity parameter
\begin{equation}
 q_g=\frac{n(n-1)}{8g\mathcal C(n/g)}.
 \label{eq:PS1}
\end{equation}
The profile endpoints in Theorem~\ref{thm:volume} imply
$q_g\asymp n^2/(g+n)\to\infty$. In particular
$q_g\sim n^2/(32\pi^2g)$ when $n/g\to0$, and
$q_g\sim x_cn/(8\pi^2)$ when $n/g\to\infty$.

We write $\operatorname{PPP}(\nu)$ for a Poisson point process with intensity measure $\nu$.

\begin{theorem}[Short-curve processes and the systole]\label{thm:pants-poisson}
\emph{(i) Curves cutting off two cusps.} The random locally finite counting measures
\begin{equation}
 \sum_{\gamma\in\mathcal P_2(X)}
       \delta_{\sqrt{q_g}\ell_\gamma}
 \ \Longrightarrow\ \operatorname{PPP}(2t\,\mathrm{d}t)
 \quad\hbox{on }[0,\infty)
 \label{eq:PS2}
\end{equation}
converge in distribution for the vague topology. Thus counts in finitely
many disjoint bounded intervals $[a_i,b_i]$ converge jointly to independent
Poisson variables with means $b_i^2-a_i^2$. If $\ell_{\min}^{(2)}$ is the
shortest length in $\mathcal P_2(X)$, with value infinity when this set is
empty, then
\[
 \mathbb P^{\mathrm{WP}}_{g,n}
 \{\sqrt{q_g}\ell_{\min}^{(2)}>t\}\longrightarrow \mathrm{e}^{-t^2}
 \qquad(t\ge0).
\]
\smallskip
\noindent\emph{(ii) All planar cusp types.} Let $\mathcal P_k(X)$ consist of simple geodesics whose
genus-zero side contains exactly $k$ cusps. If $n/g\to\lambda\in[0,\infty]$,
put $x_\lambda=X(Y_\lambda)$, where $\lambda(Y_\lambda)=\lambda$, with
$x_0=0$ and $x_\infty=x_c$. Define
\begin{equation}
 w_k(x)=\frac{2V_{0,k+1}}{k!}
             \left(\frac{x}{2\pi^2}\right)^{k-2},\qquad
 W(x)=\sum_{k\ge2}w_k(x)=\frac{2\Phi_0'(x)}{x^2}.
 \label{eq:PC1}
\end{equation}
Here $w_2(0)=W(0)=1$, $w_k(0)=0$ for $k>2$, and $W(x_c)<\infty$.
Then
\begin{equation}
 \sum_{k\ge2}\ \sum_{\gamma\in\mathcal P_k(X)}
       \delta_{(\sqrt{q_g}\ell_\gamma,k)}
 \ \Longrightarrow\
 \operatorname{PPP}\!\left(2t\,\mathrm{d}t\otimes
                    \sum_{k\ge2}w_k(x_\lambda)\delta_k\right).
 \label{eq:PC2}
\end{equation}
Convergence is for the vague topology on the product with the discrete
type space. The projection to lengths also converges, to
$\operatorname{PPP}(2W(x_\lambda)t\,\mathrm{d}t)$.

\noindent\emph{(iii) Primitive geodesics and the systole.} This is also the limit of
the process of all unoriented primitive closed geodesics, at the same
rescaling. Writing $\operatorname{sys}(X)$ for the shortest closed
geodesic length, we have
\[
 \mathbb P^{\mathrm{WP}}_{g,n}
 \{\sqrt{q_g}\operatorname{sys}(X)>t\}
 \longrightarrow \mathrm{e}^{-W(x_\lambda)t^2}\qquad(t\ge0).
\]
With probability tending to one every systolic curve is planar
separating. Select one by any measurable tie rule, and let $K_g$ be
its planar-side cusp count, with value zero if it is not planar. Then
$\bigl(\sqrt{q_g}\operatorname{sys}(X),K_g\bigr)$ converges jointly
to a pair $(Z,K)$ with
\begin{equation}
 \begin{aligned}
 \mathbb P\{Z\in \mathrm{d}t,K=k\}
 &=2t\,w_k(x_\lambda)\mathrm{e}^{-W(x_\lambda)t^2}\,\mathrm{d}t,\\
 \mathbb P\{K=k\}&=\frac{w_k(x_\lambda)}{W(x_\lambda)}
 \qquad(t>0,\ k\ge2).
 \end{aligned}
 \label{eq:FS1}
\end{equation}
\end{theorem}

The method of factorial moments for curves separating off a prescribed
number of cusps has a direct precedent in Hide--Thomas~\cite[Theorem~1.7]{HTSphere}
for punctured spheres, which also treats summed types and nested
configurations~\cite[Remark~1.9 and Proposition~3.6]{HTSphere}. Here the uniform volume ratios
give the entire family of type intensities for diverging genus at every
cusp density above the square-root scale. These planar curves
become contractible after the cusps are filled; they are excluded from
the tight spectrum of~\cite[Definition~2.1]{BL}. The all-primitive
process of Mirzakhani--Petri~\cite[Theorem~4.1]{MP} concerns closed
surfaces at fixed physical lengths. Here the full process is observed
on a shrinking scale; a tight-spectrum limit at its longer natural
scale is a separate question.

\emph{Factorial moments.} For sufficiently small $L$, the collar lemma forces all curves counted by
$N_L=\#\{\gamma\in\mathcal P_2(X):\ell_\gamma\le L\}$ to be disjoint.
Their pants interiors are also disjoint. Otherwise one pants boundary
would be an interior essential geodesic of the other pants, hence
boundary-parallel, contradicting uniqueness of geodesic representatives.
Thus their pairs of cusp labels are disjoint, and their complement has
type $(g,n-r)$ for $r$ pants. We have $g\ge2$ eventually, so the
exceptional double-pants decomposition of a four-punctured sphere does
not arise.

Mirzakhani's integration formula~\cite{MirInv,Mir}, with the same conventions as
\cite[Lemma~29]{SW}, gives, for any fixed $r$ and short length windows
$I_1,\ldots,I_r$, the ordered distinct-curve integral
\begin{equation}
 \mathbb E\!\sum_{\substack{\gamma_1,\ldots,\gamma_r\in\mathcal P_2(X)\\
                            \text{distinct}}}
       \prod_{j=1}^r\mathbf1_{\ell_{\gamma_j}\in I_j}
 =\frac{(n)_{2r}}{2^r V_{g,n}}
    \int_{I_1\times\cdots\times I_r}
       V_{g,n-r}(x_1,\ldots,x_r,0^{n-2r})
       \prod_{j=1}^r x_j\,\mathrm{d}x_j.
 \label{eq:PS3}
\end{equation}
Here $(n)_{2r}=n(n-1)\cdots(n-2r+1)$. The factor $2^{-r}$ removes
the order of the two labels within each pair. The tuple of curves is
ordered, so there is no $r!$ divisor. Each pants volume is $V_{0,3}=1$.
Mirzakhani's boundary estimate~\cite[Lemma~3.2]{Mir}, in the form
of~\cite[Lemma~14]{SW},
\begin{equation}
 1\le\frac{V_{g,m}(\mathbf x)}{V_{g,m}}
       \le\exp\!\left(\frac12\sum_jx_j\right)
 \label{eq:PS4}
\end{equation}
is uniform in the topological type. For fixed $r$, it replaces the
boundary volume in \eqref{eq:PS3} by $V_{g,n-r}(1+O_r(L))$ when
all $I_j\subset[0,L]$.

Theorem~\ref{thm:volume} also implies, for every fixed $r$,
\begin{equation}
 \frac{V_{g,n-r}}{V_{g,n}}
  =(2g\mathcal C(n/g))^{-r}(1+o(1)).
 \label{eq:PS5}
\end{equation}
Indeed this is a product of $r$ adjacent cusp ratios. Replacing $(n-j)/g$
by $n/g$ changes their profile by $1+O_r(n^{-1})$, because
$\lambda\partial_\lambda\log\mathcal C(\lambda)$ is bounded: this follows
from analyticity in the interior and the two endpoint expansions.

\begin{proof}[Proof of Theorem~\ref{thm:pants-poisson}]
We prove the planar process by bounding the tail over cusp types and
computing the mixed factorial moments for any fixed collection of types.
We then exclude the other topological types on the same length scale;
this identifies the full primitive process and the systole. Write $P=2\pi^2$, $h=2g-1$, $m=n-1$, and $F_h=\Phi_g'$.
Choose $x\in(0,x_c)$ by $h\Psi(x)=m$ and put $q_g^*=nx/(4P)$.
The effective density is $2(n-1)/(2g-1)$. The boundedness of
$\lambda\partial_\lambda\log X(Y_\lambda)$, including its endpoint limits, implies
$q_g^*/q_g\to1$ and $q_g^*\to\infty$.

\emph{Step 1: a uniform tail bound over cusp types.}
We first show that curves cutting off a large number of cusps contribute
negligibly, uniformly along the sequence. The input is the
coefficient-ratio bound \eqref{eq:PC3}, which compares volumes whose cusp
counts differ by $j$.
By the positive convolution \eqref{eq:TL33} and
Lemma~\ref{lem:point-mass}, every point mass is at most $C/\sigma$,
where $\sigma^2=h\mathscr D\Psi(x)$. The local limit
\eqref{eq:TL47} gives a lower bound $c/\sigma$ at $m$.
Consequently, for all sufficiently large $h,m$,
\begin{equation}
 \frac{[z^{m-j}]F_h(z)}{[z^m]F_h(z)}
 =x^j\frac{\mathbb P_x(N=m-j)}{\mathbb P_x(N=m)}
 \le Cx^j\qquad(0\le j\le m).
 \label{eq:PC3}
\end{equation}
Only the denominator uses the local limit; the numerator may be at any
index, including zero.

Let $N_{k,L}$ count curves of type $k$ with length at most $L$, and put
$a_k=V_{0,k+1}/(k!P^{k-2})$. Mirzakhani integration gives
\begin{equation}
 \begin{aligned}
 \mathbb E N_{k,L}=\frac{\binom nk}{V_{g,n}}
 \int_0^L\ell\,V_{0,k+1}(\ell,0^k)
                  V_{g,n-k+1}(\ell,0^{n-k})\,\mathrm{d}\ell,\\
 \binom nk\frac{V_{0,k+1}V_{g,n-k+1}}{V_{g,n}}
 =\frac nP a_k\frac{[z^{n-k}]F_h}{[z^{n-1}]F_h}.
 \end{aligned}
 \label{eq:PC4}
\end{equation}
The two sides have different genera, so there is no side-exchange
factor. The coefficient identity follows from \eqref{eq:TL50}; it
holds up to $k=n$. Bound \eqref{eq:PS4} changes the integrand's two
volumes by a factor between one and $\mathrm{e}^L$, independently of $k$.
Using \eqref{eq:PC3} with $j=k-1$ therefore gives, for each fixed $T$,
\begin{equation}
 \mathbb E\sum_{k\ge K}N_{k,T/\sqrt{q_g^*}}
 \le CT^2\sum_{k\ge K}a_kx_c^{k-2}\longrightarrow0
 \quad(K\to\infty),
 \label{eq:PC5}
\end{equation}
uniformly along the sequence. To justify summability directly, the
genus-zero string solution gives $\Phi_0'''(x)=1/a(Y(x))=Y'(x)$.
Since $V_{0,3}=1$, its zero initial conditions give
$\Phi_0''=Y$ and $\Phi_0'(x)=\int_0^xY(u)\,\mathrm{d}u$. This is finite at
$x_c$; positivity gives convergence of the coefficient series there.
Its initial term $x^2/2$ also gives $W(0)=1$.

\emph{Step 2: fixed types and nested configurations.}
Next we compute the mixed factorial moments for finitely many fixed cusp
types. The main contribution comes from configurations whose planar
sides are pairwise disjoint, which we call \emph{stars}; nested
configurations are shown to be negligible.
Fix $r$ ordered distinct curves of prescribed types $k_1,\ldots,k_r$,
and write $K=\sum k_i$. The collar lemma makes their boundaries
disjoint. For a star, the ordered
cusp-label multiplicity is $(n)_K/\prod_i k_i!$, and the complement
has type $(g,n-K+r)$. There is no $r!$ divisor. The boundary error is
$1+O_r(L)$, as each cut occurs twice in \eqref{eq:PS4}.
For fixed $d$, \eqref{eq:PS5} and $q_g^*/q_g\to1$ give
$V_{g,n-d}/V_{g,n}=(x/(Pn))^d(1+o(1))$.
Thus for windows $[a_i,b_i]/\sqrt{q_g^*}$ the star contribution tends to
\begin{equation}
 \prod_{i=1}^r
       w_{k_i}(x_\lambda)(b_i^2-a_i^2).
 \label{eq:PC6}
\end{equation}
The positive expression has a relative error before taking the limit,
even when some limiting weights vanish.

Filling cusps turns each planar side into a disk. Disjoint boundaries
therefore have nested or disjoint disks; their cut graph is a tree with
one genus-$g$ root and $r$ planar vertices. Let $d$ be the root degree,
$a_v$ the number of cusps at a planar vertex and $d_v$ its degree;
put $a=\sum_v a_v$. Prescribed types bound $a$ by $K$.
There are finitely many such trees with ordered edges and cusp
allocations, depending only on the fixed types. Classification of
surfaces gives the corresponding mapping-class orbits with bounded
symmetry factors; assigning cusp labels costs at most $Cn^a$.
The root volume ratio is $O((x/(Pn))^{a-d})$, the planar volumes
are fixed, and the $r$ length integrals cost $O((nx)^{-r})$.
Each cut type contributes at most
\begin{equation}
 Cn^{d-r}x^{a-d-r}=O(n^{-1})\quad\hbox{if it is not a star}.
 \label{eq:PC7}
\end{equation}
Indeed stability gives $a_v+d_v\ge3$ and
$\sum_vd_v=2r-d$, hence $a\ge r+d$. The power of $x$ is
nonnegative, so no loss occurs as $x\to0$. A nonstar has $d\le r-1$;
if $d=r$, every planar vertex is a leaf and the configuration is a star.
This accounts for every fixed-type mixed factorial moment.

The limits \eqref{eq:PC6} are the mixed factorial moments of independent
Poisson counts. Higher fixed moments give uniform integrability, and
the product Poisson law is moment determinate. First moments give
tightness on compact sets in the marked space. This proves
\eqref{eq:PC2}. Its unmarked projection requires more than vague
convergence: truncate the types, pass to the finite-type limit, then
let the cutoff grow using \eqref{eq:PC5} and the convergent limiting
series. This proves the stated unmarked process. Finally replace
$q_g^*$ by $q_g$. The type-two restriction gives \eqref{eq:PS2}
without assuming a limit for $n/g$, by subsequences and compactness of
$[0,\infty]$. Its zero-count probabilities give the pants minimum law.

\emph{Step 3: the remaining curves and the systole.}
It remains to show that, on the scale $q_g^{-1/2}$, there are no short
curves other than planar separating ones: nonseparating curves, and
separating curves with positive genus on both sides, are rare.
Nonseparating simple curves of length at most $L\le L_0$ have expected
count at most $C\mathrm{e}^LL^2V_{g-1,n+2}/V_{g,n}\le CL^2$, by integration,
\eqref{eq:PS4}, and the uniform lower bound for $B_{g,n}$.
To sum every positive-genus separating cut, put
$S_g(z)=\sum_{i=1}^{g-1}\Phi_i'(z)\Phi_{g-i}'(z)$. Exact positivity
\eqref{eq:TL33} includes genus one, so the child powers add to $h-1$:
\begin{equation}
 S_g(z)=V(z)^{h-1}Q_g(z),\qquad Q_g\ge0
 \ \hbox{coefficientwise},\qquad
 \frac{S_g(x)}{F_h(x)}\le\frac{Ca(Y)r(Y)^2}{g^2}.
 \label{eq:FS2}
\end{equation}
The last bound is \eqref{eq:TL26}: undoing \eqref{eq:TL9} leaves
exactly $ar^2$, as the child complexities sum to $2g-2$ and their
deficits to $3g-4$. In particular bounded child genera are retained.
Lemma~\ref{lem:point-mass}, applied to the positive factor $V^{h-1}$,
gives $\sup_j\mathbb P_x^{S}(N=j)\le C/\sigma$. Against the parent lower
bound $\mathbb P_x^{F}(N=n-1)\ge c/\sigma$, the standard deviations
cancel. The coefficient identity \eqref{eq:TL50} therefore yields
\begin{equation}
 \begin{aligned}
 \mathcal A_{g,n}
 &:=\sum_{i=1}^{g-1}\sum_{k=0}^n\binom nk
       \frac{V_{i,k+1}V_{g-i,n-k+1}}{V_{g,n}}
   =\frac1P\frac{[z^n]S_g(z)}{[z^n]\Phi_g(z)}\\
 &=\frac{n}{Px}\frac{S_g(x)}{F_h(x)}
       \frac{\mathbb P_x^S(N=n)}{\mathbb P_x^F(N=n-1)}
   \le\frac{Cnar^2}{xg^2}\le\frac Cg.
 \end{aligned}
 \label{eq:FS3}
\end{equation}
Indeed $(n-1)/x=h/(Aar)$, while $A$ is bounded below and $r$ above.
Thus no small-density factor is lost. Each separating orbit with two
positive genera occurs in this ordered sum. Its integration symmetry
and one-handle factors are at most one; applying \eqref{eq:PS4} to
both components bounds its total expected count by
$C\mathrm{e}^LL^2\mathcal A_{g,n}\le CL^2/g$. This includes all cusp partitions.

For fixed $T$, take $L=T/\sqrt{q_g}$. Markov's inequality shows that
with probability $1-O(T^2/q_g)$ every simple geodesic in this window
is planar. Primitive geodesics shorter than $2\operatorname{arcsinh}1$
are simple, by the standard collar theorem used in
\cite[Section~1.1]{HTSphere}. Since $q_g\to\infty$, the full primitive
process and planar process agree on each bounded rescaled window
with probability tending to one. This transfers the process limit
and its zero-count law to the systole.

Finally restrict the first point to $[\eta,T]$ and finitely many cusp
types. The first-point map is continuous when its first length is
unique and avoids the endpoints, as happens almost surely in the
limiting Poisson process. Let the type cutoff grow using
\eqref{eq:PC5}, then $\eta\downarrow0$ and $T\to\infty$.
The omitted length probabilities tend to zero by the zero-count law;
$W\ge1$. Nonplanar competitors are excluded on each fixed window.
This also makes any tie rule immaterial and proves planarity of every
systolic curve with probability tending to one. The first-arrival
density of the limiting marked process is \eqref{eq:FS1}.
\end{proof}

\subsection{Many small eigenvalues}\label{sec:pants-eigenvalues}

\begin{proposition}[Many short pants]\label{prop:pants-concentration}
If $L=L_g\to0$ and $q_gL_g^2\to\infty$, then
\begin{equation}
 \frac{N_{L_g}}{q_gL_g^2}\longrightarrow1
 \quad\hbox{in probability}.
 \label{eq:PS6}
\end{equation}
Consequently, for each fixed $\varepsilon\in(0,1)$, with probability
tending to one there are at least
$\lfloor(1-\varepsilon)q_gL_g^2\rfloor-1$ positive discrete eigenvalues
in $(0,6L_g]$, counted with multiplicity.
\end{proposition}
\begin{proof}
The first two cases of \eqref{eq:PS3} give
\[
 m_L:=\mathbb E N_L=\frac{(n)_2L^2}{4}
       \frac{V_{g,n-1}}{V_{g,n}}(1+O(L))\sim q_gL^2,
\]
and
\[
 \frac{\mathbb E(N_L)_2}{m_L^2}
 =\frac{(n-2)(n-3)}{n(n-1)}D_{g,n-1}^{-1}(1+O(L))\longrightarrow1.
\]
Since $m_L\to\infty$, the relative variance tends to zero, proving
\eqref{eq:PS6}. This second-moment argument does not use the point-process
limit at a window growing with $g$.

Hide--Thomas~\cite[Lemma~5.3]{HTSphere} prove that $N$ such pants
boundaries of length at most $L$ give $N-1$ positive eigenvalues below
$6L$, provided $6L<1/4$. For the reader's convenience we recall the
variational argument.
Inside each pants choose a function equal to inward distance $\rho$ on
the half-collar $0\le\rho\le1$, equal to one beyond it, and zero outside
the pants. For all sufficiently small boundary lengths $\ell$, the
collar lemma supplies this embedded half-collar. The function lies in
$H^1(X)$, and
\[
 \int|\nabla f|^2=\ell\sinh1,\qquad
 \int|f|^2\ge2\pi-\ell\sinh1.
\]
Different pants have disjoint supports up to their boundaries, so the
span of these functions has dimension $N$ and Rayleigh quotient at most
$6L$ for small $L$. Its intersection with the orthogonal complement of
constants has dimension at least $N-1$. Since the essential spectrum
starts at $1/4$, the form version of min--max gives the claimed positive
discrete eigenvalues. Apply \eqref{eq:PS6} and let $g$ grow.
\end{proof}

\begin{corollary}[A growing number of small eigenvalues]\label{cor:many-spectrum}
Suppose $g\to\infty$ and $n/\sqrt g\to\infty$. If $j=j(g)\to\infty$
and $j=o(q_g)$, then
\[
 \mathbb P^{\mathrm{WP}}_{g,n}
 \{\lambda_j(X)\text{ exists and }\lambda_j(X)\le12\sqrt{j/q_g}\}
 \longrightarrow1,
\]
where positive discrete eigenvalues are ordered with multiplicity.
\end{corollary}

\begin{proof}[Proof of Corollary~\ref{cor:many-spectrum}]
For $j\to\infty$ with $j=o(q_g)$, choose $L_g=2\sqrt{j/q_g}$.
Then $L_g\to0$, $q_gL_g^2=4j$, and \eqref{eq:PS6} implies
$N_{L_g}\ge j+1$ with probability tending to one. The preceding
variational estimate gives $\lambda_j\le6L_g=12\sqrt{j/q_g}$.
\end{proof}

The spectral mechanism is established in the fixed-genus work of
Hide--Thomas~\cite[Theorem~1.10 and Remark~5.5]{HTSphere}; their later
bordered-volume asymptotics give linear counts below each fixed positive
threshold at fixed genus~\cite[Theorem~1.1]{HTLargeN}.
Their deterministic result~\cite[Theorem~1.3]{HTMany}, in the July 2026
revision, gives a linear number below $1/4$ whenever $g<an$ for fixed
$a>0$. Our threshold tends to zero and the genus grows. In the range
$\sqrt g\ll n\ll g$, the corresponding second-moment volume inputs
already follow from Shen--Wu~\cite{SW}. The extension supplied here is
the density-dependent intensity and removal of the upper restriction on
$n/g$. No sharp eigenvalue asymptotic, matching lower bound, or
Poisson eigenvalue law is asserted.


\subsection{The tight length spectrum}\label{subsec:tight-spectrum}

We now consider curves that remain noncontractible after filling all
cusps. The corresponding tight Weil--Petersson volumes were introduced
by Budd~\cite{Budd}. Following Budd and Lions~\cite{BL}, a closed geodesic is
\emph{tight} if it is the unique shortest representative of its free
homotopy class on the filled surface. The distinction matters: several
ordinary geodesics can belong to the same class after the cusps are
filled. The tight integration formula and the geometric argument used
in the proof are due to~\cite{BL}; the new input is a volume estimate that
allows arbitrary prescribed cusp counts.

Let $Y_\lambda$ satisfy $\lambda(Y_\lambda)=\lambda$,
as in Proposition~\ref{prop:bounded-boundaries}, and set
\begin{equation}
 \alpha(\lambda)=\frac{\sqrt{2r(Y_\lambda)}}\pi,
 \qquad \nu(\mathrm dt)=\frac{\cosh t-1}{t}\,\mathrm dt.
 \label{eq:tight-scale}
\end{equation}
Here $\alpha(0)=1$, and the density of $\nu$ is defined to be zero
at the origin.

\begin{theorem}[Tight lengths at arbitrary cusp density]
\label{thm:tight-spectrum}
Let $n=n(g)\ge0$ be any integer sequence, and let $X$ have law
$\mathbb P^{\mathrm{WP}}_{g,n}$. As $g\to\infty$, the point process
\begin{equation}
 \sum_{\substack{\gamma\text{ primitive, unoriented}\\
                  \gamma\text{ tight on }X}}
 \delta_{\alpha(n/g)\ell_X(\gamma)}
 \ \Longrightarrow\ \operatorname{PPP}(\nu)
 \label{eq:tight-process}
\end{equation}
converges in distribution for the vague topology on $[0,\infty)$.
For every fixed $L<\infty$, with probability tending to one all the
curves counted in $[0,L]$ are simple, pairwise disjoint, and have
connected complement after filling the cusps.
\end{theorem}

In particular, no limit of $n/g$ is required. At large density,
the endpoint expansion of $r$ gives
\[
 \alpha(\lambda)\sim
 \frac{\sqrt{3j_{0,1}\sqrt5}}\pi\lambda^{-1/4}.
\]
This agrees with the scale in~\cite[Corollary~1.3]{BL}.
Budd and Lions asked whether their limit persists for other cusp
weights and under the fixed-count measures
$\mathbb P^{\mathrm{WP}}_{g,n}$~\cite[Open Question~1.6]{BL}.
Theorem~\ref{thm:tight-spectrum} treats the fixed-count part of this
question; Corollary~\ref{cor:tight-boltzmann} below treats the
remaining cusp-weight regimes. The planar curves in
Theorem~\ref{thm:pants-poisson} are contractible after filling the
cusps and therefore do not occur in \eqref{eq:tight-process}.

The proof follows the moment method of Budd and Lions. Their tight
integration formula~\cite[Proposition~3.2]{BL} expresses factorial moments
of tight curve counts through tight volumes of the cut surfaces. We first
write the generating functions of tight volumes through the shifted
amplitudes of Section~\ref{exact-amplitudes-and-the-harmonic-model}
(Lemma~\ref{lem:tight-gf}); then extract the coefficient at the
prescribed cusp count, as in Section~\ref{extracting-fixed-cusp-counts}
(Lemma~\ref{lem:tight-connected}); then show that configurations with
disconnected complement are negligible (Lemma~\ref{lem:tight-separating});
and finally identify the Poisson limit (Subsection~\ref{subsec:tight-proof}).
These steps are carried out in Appendix~\ref{app:tight}.

\emph{Random cusp counts.}
For $0\le\mu<\mu_{\rm cr}=x_c/(2\pi^2)$, the Boltzmann measure
$\mathbb P^{\mathrm{WP}}_{g,\mu}$ first chooses the cusp count $N$
with probabilities proportional to $\mu^nV_{g,n}/n!$, then samples
$\mathbb P^{\mathrm{WP}}_{g,n}$ conditionally on $N=n$.
Its probability generating function is
$\Phi_g(2\pi^2\mu z)/\Phi_g(2\pi^2\mu)$, as in~\cite{BL}.

\begin{corollary}[Arbitrary subcritical cusp weights]
\label{cor:tight-boltzmann}
For any sequence $0\le\mu_g<\mu_{\rm cr}$, put
\[
 a_g=\frac{\sqrt{2r(Y(2\pi^2\mu_g))}}\pi.
\]
Under $\mathbb P^{\mathrm{WP}}_{g,\mu_g}$, the primitive tight
lengths multiplied by $a_g$ converge to $\operatorname{PPP}(\nu)$.
In particular, if $\mu_g\to\mu<\mu_{\rm cr}$, the unscaled lengths have
limiting intensity $(\cosh(a_\mu t)-1)\,\mathrm dt/t$, where
$a_\mu=\sqrt{2r(Y(2\pi^2\mu))}/\pi$. If $\mu_g\to\mu_{\rm cr}$, no condition on
the speed of convergence is required.
\end{corollary}
\noindent The proof is given in Subsection~\ref{subsec:tight-boltzmann-proof}.


\subsection{The full length spectrum at the critical cusp scale}
\label{sec:critical-cluster}

The tight spectrum omits curves that become contractible when the cusps
are filled. At $n\asymp\sqrt g$, these curves contribute on the same
length scale as the nonseparating curves. Their contribution is not
Poisson: a two-cusp pair of pants contributes its whole length spectrum,
not only its boundary.

For $s>0$, let $P_s$ be the hyperbolic pair of pants with two cusps and
one geodesic boundary of length $s$. Write $K_s$ for its primitive
unoriented length counting measure, including the boundary once and
excluding parabolic classes. Equal lengths are counted with their
geometric multiplicities. Let
\begin{equation}
 \mu_c(\mathrm ds)=\frac{c^2}{8\pi^2}\sinh(s/2)\,\mathrm ds,
 \qquad \nu(\mathrm dt)=\frac{\cosh t-1}{t}\,\mathrm dt.
 \label{eq:cluster-intensities}
\end{equation}

\begin{theorem}[Critical primitive length spectrum]
\label{thm:critical-cluster}
Suppose $g\to\infty$ and $n/\sqrt g\to c\in(0,\infty)$, and let
$X$ have law $\mathbb P^{\mathrm{WP}}_{g,n}$. Then
\begin{equation}
 \sum_{\substack{\gamma\text{ primitive}\\\text{unoriented on }X}}
       \delta_{\ell_X(\gamma)}
 \ \Longrightarrow\ \Xi_c:=\mathcal N+\sum_{s\in\mathcal Q_c}K_s,
 \label{eq:critical-cluster}
\end{equation}
as locally finite counting measures on $[0,\infty)$ with the vague
topology. Here $\mathcal N$ and $\mathcal Q_c$ are independent Poisson
point processes of intensities $\nu$ and $\mu_c$. In particular, for
every nonnegative continuous compactly supported function $f$,
\begin{equation}
 \mathbb E\mathrm e^{-\langle f,\Xi_c\rangle}
 =\exp\left\{-\int(1-\mathrm e^{-f(t)})\,\nu(\mathrm dt)
 -\int(1-\mathrm e^{-\langle f,K_s\rangle})\,\mu_c(\mathrm ds)\right\}.
 \label{eq:cluster-laplace}
\end{equation}
\end{theorem}

Shen and Wu obtained the factorial-moment asymptotics for short
two-cusp pants and the limiting probability that such a curve occurs
at this cusp scale~\cite[Lemmas~28 and~33]{SW}. The corresponding Poisson mean is
$\mu_c([0,L])=c^2(\cosh(L/2)-1)/(4\pi^2)$ on their stated short
window. Theorem~\ref{thm:critical-cluster} concerns the full primitive
process on every fixed length window, including the correlated lengths
inside each pair of pants. The background intensity and the geometric
moment method come from Mirzakhani and Petri~\cite{MP}. The critical
regime needs the volume bounds of Mirzakhani--Zograf and
Shen--Wu~\cite{MZ,SW} together with the zero-density case of
Theorem~\ref{thm:volume} and Proposition~\ref{prop:bounded-boundaries}:
at $n\asymp\sqrt g$ the earlier estimates bound the genus ratio
$V_{g-1,n+2}/V_{g,n}$ only from above. The full uniform expansion is not
used.

The nonsimple part distinguishes this regime from closed surfaces.
He, Shen, Wu and Xue proved that the nonsimple systole of a
Weil--Petersson random closed surface is asymptotic to $\log g$,
with a more precise concentration about $\log g-\log\log g$~\cite[Theorem~1.1]{HSWX}.
Thus nonsimple curves leave every
fixed length window in that model. At the critical cusp scale here,
two-cusp pants instead retain nonsimple curves of bounded length.

\emph{Outline of the proof.} Short multicurves are organized by their cut
graphs. Using only the volume bounds of~\cite{MZ,SW},
Lemma~\ref{lem:critical-cuts} shows that in the limit a cut graph
contributes only if it has one component carrying the genus, possibly
with nonseparating loops, and two-cusp pairs of pants attached to it as
leaves. A bounded-multiplicity argument (Lemmas~\ref{lem:primitive-count}
and~\ref{lem:critical-complex}) reduces intersecting simple curves and
nonsimple curves to such cut graphs; hence, with probability tending to
one, every short nonsimple curve lies in a two-cusp pair of pants. The mixed
factorial moments of nonseparating lengths and pants waists converge to those
of independent Poisson processes (Lemma~\ref{lem:critical-skeleton}, where the
zero-density volume ratios enter), and attaching to each waist the full
spectrum of its pair of pants gives $\Xi_c$. The details are given in
Appendix~\ref{app:critical}.

The limit is not Poisson. If $m_L(s)=K_s([0,L])$, then
Lemma~\ref{lem:primitive-count} and the support restriction $s\le2L$
give finite moments and
\begin{equation}
 \operatorname{Var}\Xi_c([0,L])-\mathbb E\Xi_c([0,L])
 =\int m_L(s)(m_L(s)-1)\,\mu_c(\mathrm ds).
 \label{eq:cluster-overdispersion}
\end{equation}
Besides its waist, $P_s$ contains the primitive figure-eight around
its two cusps, of length $2\operatorname{arccosh}(2+\cosh(s/2))$
by \cite[Theorem~1.3]{Rivin}. For $L>2\operatorname{arccosh}3$,
both lengths lie below $L$ on an interval of positive $\mu_c$-measure.
The right side of \eqref{eq:cluster-overdispersion} is therefore
strictly positive. The theorem identifies these correlations in the
limit; no convergence of the full-spectrum second moments is asserted.

